% =============================================================================
%  Section 3 — Direct Repository Access: Claude Sonnet 4.6 (Raw Run)
%  Drop into the full report alongside section3_raw_runs.tex (Opus raw)
% =============================================================================

\section{Direct Repository Access Evaluation --- Claude Sonnet~4.6 (Raw)}\label{sec:raw-sonnet}

\subsection{Methodology}

Each of the 22 \texttt{astropy/astropy} tasks from the
\texttt{princeton-nlp/SWE-bench\_Verified} dataset was presented to
Claude Code (model: \texttt{claude-sonnet-4-6}) with \emph{direct
filesystem access} to a locally cloned copy of the repository, checked
out at the exact base commit specified by the task.
No MCP knowledge graph, no web search, no \texttt{git} CLI, and no
pre-computed answer cache were available to the model.

\paragraph{Prompt structure.}
Every task was issued under a fixed set of ground rules:

\begin{enumerate}
  \item \textbf{Read files directly} from the local clone at the
        base-commit path.
  \item \textbf{Form the answer entirely from first-principles evidence}
        gathered during the session.
  \item \textbf{No web search}, \textbf{no \texttt{git} commands}
        (\texttt{git log}, \texttt{git blame}, \texttt{git diff}, etc.),
        and \textbf{no remote API calls}.
  \item \textbf{No reading} from \texttt{answers/}, \texttt{results/},
        \texttt{expected/}, or \texttt{cached/} directories.
\end{enumerate}

\noindent
The SWE-Bench question text (issue description, hints, and required
failing tests) was appended after the ground rules.
Claude Code was then free to explore the codebase, read any source file,
and produce a unified diff patch as its answer, written to
\texttt{answer.json} alongside telemetry (time, cost, per-model token
counts).

\paragraph{Judging.}
Answers were graded against the ground-truth patch on the same
three-dimensional rubric used across all runs
(Table~\ref{tab:rubric-raw-sonnet}), scored by a separate Claude
Sonnet~4.6 judge session with access to \texttt{astropy\_tasks.json}.

\begin{table}[h]
\centering
\caption{Scoring rubric --- 8-point scale.}\label{tab:rubric-raw-sonnet}
\begin{tabular}{lcp{8cm}}
\toprule
\textbf{Dimension} & \textbf{Max} & \textbf{Criterion} \\
\midrule
Root cause identification   & 3 & Correctly diagnoses \emph{why} the bug exists. \\
Correct file(s)             & 2 & Identifies the right file(s) to change. \\
Correct patch / code change & 3 & Proposed change matches ground truth or is functionally equivalent. \\
\bottomrule
\end{tabular}
\end{table}

\noindent
Grade tiers: \textbf{Exact}~(8/8), \textbf{Near-perfect}~(7/8),
\textbf{Partial}~(5--6/8), \textbf{Fail}~($\le$4/8).

% ---------------------------------------------------------------------------
\subsection{Results}

\subsubsection{Aggregate Performance}

\begin{table}[h]
\centering
\caption{Dimension-level scores across all 22 Sonnet raw tasks.}\label{tab:dim-raw-sonnet}
\begin{tabular}{lrrr}
\toprule
\textbf{Dimension} & \textbf{Score} & \textbf{Max} & \textbf{\%} \\
\midrule
Root cause identification   & 66 & 66 & 100\%  \\
Correct file(s)             & 44 & 44 & 100\%  \\
Correct patch / code change & 54 & 66 & 81.8\% \\
\midrule
\textbf{Overall}            & \textbf{164} & \textbf{176} & \textbf{93.2\%} \\
\bottomrule
\end{tabular}
\end{table}

\begin{table}[h]
\centering
\caption{Grade distribution (Sonnet raw).}\label{tab:grade-raw-sonnet}
\begin{tabular}{lcl}
\toprule
\textbf{Grade} & \textbf{Count} & \textbf{Tasks} \\
\midrule
Exact (8/8)        & 14 & 1, 3--7, 9, 12, 14--16, 18, 20, 22 \\
Near-perfect (7/8) &  5 & 2, 10, 11, 17, 21 \\
Partial (6/8)      &  2 & 8, 13 \\
Partial (5/8)      &  1 & 19 \\
Fail ($\le$4/8)    &  0 & --- \\
\bottomrule
\end{tabular}
\end{table}

\subsubsection{Per-Task Scores, Time \& Cost}

\begin{table}[h]
\centering
\small
\caption{Per-task Sonnet raw results --- score, wall-clock time, and API cost.}\label{tab:per-task-raw-sonnet}
\setlength{\tabcolsep}{4pt}
\begin{tabular}{clccccrr}
\toprule
\textbf{\#} & \textbf{Instance ID} & \textbf{RC} & \textbf{Files} & \textbf{Patch} & \textbf{Score} & \textbf{Time (s)} & \textbf{Cost (\$)} \\
\midrule
 1 & \texttt{astropy-12907} & 3 & 2 & 3 & 8/8 &     87 & 0.187 \\
 2 & \texttt{astropy-13033} & 3 & 2 & 2 & 7/8 &    372 & 0.601 \\
 3 & \texttt{astropy-13236} & 3 & 2 & 3 & 8/8 &    236 & 0.577 \\
 4 & \texttt{astropy-13398} & 3 & 2 & 3 & 8/8 &    881 & 1.849 \\
 5 & \texttt{astropy-13453} & 3 & 2 & 3 & 8/8 &     98 & 0.227 \\
 6 & \texttt{astropy-13579} & 3 & 2 & 3 & 8/8 &    204 & 0.374 \\
 7 & \texttt{astropy-13977} & 3 & 2 & 3 & 8/8 &    247 & 0.531 \\
 8 & \texttt{astropy-14096} & 3 & 2 & 1 & 6/8 &    209 & 0.312 \\
 9 & \texttt{astropy-14182} & 3 & 2 & 3 & 8/8 &    371 & 0.654 \\
10 & \texttt{astropy-14309} & 3 & 2 & 2 & 7/8 &    130 & 0.156 \\
11 & \texttt{astropy-14365} & 3 & 2 & 2 & 7/8 &    168 & 0.273 \\
12 & \texttt{astropy-14369} & 3 & 2 & 3 & 8/8 &    836 & 1.300 \\
13 & \texttt{astropy-14508} & 3 & 2 & 1 & 6/8 &    141 & 0.262 \\
14 & \texttt{astropy-14539} & 3 & 2 & 3 & 8/8 &     22 & 0.053 \\
15 & \texttt{astropy-14598} & 3 & 2 & 3 & 8/8 & 1{,}637 & 2.381 \\
16 & \texttt{astropy-14995} & 3 & 2 & 3 & 8/8 &     53 & 0.128 \\
17 & \texttt{astropy-7166}  & 3 & 2 & 2 & 7/8 &    112 & 0.212 \\
18 & \texttt{astropy-7336}  & 3 & 2 & 3 & 8/8 &     55 & 0.104 \\
19 & \texttt{astropy-7606}  & 3 & 2 & 0 & 5/8 &    100 & 0.183 \\
20 & \texttt{astropy-7671}  & 3 & 2 & 3 & 8/8 &    512 & 0.939 \\
21 & \texttt{astropy-8707}  & 3 & 2 & 2 & 7/8 &    261 & 0.371 \\
22 & \texttt{astropy-8872}  & 3 & 2 & 3 & 8/8 &     94 & 0.155 \\
\midrule
   & \textbf{Total}  & \textbf{66/66} & \textbf{44/44} & \textbf{54/66} & \textbf{164/176} & \textbf{6{,}826} & \textbf{11.83} \\
   & \textbf{Average} &       &       &       & \textbf{7.5/8}   & \textbf{310} & \textbf{0.54} \\
\bottomrule
\end{tabular}
\end{table}

% ---------------------------------------------------------------------------
\subsubsection{Token \& Cost Breakdown (Aggregated)}

\begin{table}[h]
\centering
\small
\caption{Aggregated token usage and cost across all 22 Sonnet raw tasks.}\label{tab:tokens-agg-raw-sonnet}
\setlength{\tabcolsep}{4pt}
\begin{tabular}{lrrrrr}
\toprule
& \textbf{Input} & \textbf{Cache Write} & \textbf{Cache Read} & \textbf{Output} & \textbf{Cost (\$)} \\
\midrule
Total   &  1{,}806 & 1{,}015{,}628 & 14{,}963{,}747 & 432{,}443 & 11.83 \\
Average &       82 &    46{,}165   &    680{,}170   &  19{,}656 &  0.54 \\
\bottomrule
\end{tabular}
\\[4pt]
{\footnotesize
Effective Weighted Input = Input + $1.25 \times$ Cache Write + $0.1 \times$ Cache Read.\quad
Effective Input = Input + Cache Write + Cache Read.}
\end{table}

% ---------------------------------------------------------------------------
\subsubsection{Per-Model Token Usage}

As in the Opus raw run, Claude Code dispatches requests internally to
\texttt{claude-sonnet-4-6} for reasoning and patch generation, and
\texttt{claude-haiku-4-5} for lightweight triage and tool routing.
Unlike the Opus raw run where Haiku handled exactly one request on
11 of 22 tasks, here Haiku is invoked only on tasks with larger
codebase surface area; on 10 of 22 tasks the run is entirely
Sonnet-only (Table~\ref{tab:model-split-raw-sonnet}).

\begin{table}[h]
\centering
\caption{Token and cost split between Sonnet and Haiku (raw run).}\label{tab:model-split-raw-sonnet}
\begin{tabular}{lrrrr}
\toprule
\textbf{Metric} & \textbf{Sonnet} & \textbf{Haiku} & \textbf{Total} & \textbf{Sonnet \%} \\
\midrule
Input tokens       &      1{,}064 &          742 &      1{,}806 &  58.9\% \\
Cache Write tokens &    606{,}905 &    407{,}723 &  1{,}015{,}628 &  59.7\% \\
Cache Read tokens  &  8{,}136{,}814 & 6{,}826{,}933 & 14{,}963{,}747 &  54.4\% \\
Output tokens      &    375{,}034 &     57{,}410 &    432{,}443 &  86.7\% \\
Cost (USD)         &     \$10.351 &      \$1.480 &     \$11.830 &  87.5\% \\
API requests       &          316 &          210 &          526 &  60.1\% \\
\bottomrule
\end{tabular}
\end{table}

\noindent
Sonnet generates 86.7\% of all output tokens and accounts for 87.5\% of
total cost.  In contrast to the Opus raw run where Opus claimed 92.1\%
of cost, Haiku's share is larger here (12.5\% vs.\ 7.9\%) because the
Sonnet tier's lower per-token rate narrows the absolute cost gap between
the two models.

% ---------------------------------------------------------------------------
\subsection{Observations}\label{sec:raw-sonnet-obs}

The following observations are derived strictly from the Sonnet~4.6 raw
run results and, where noted, from comparison with the Opus~4.6 raw run
(Section~\ref{sec:raw}).

\begin{enumerate}

\item \textbf{Highest overall accuracy of any raw run.}
      Sonnet scored 164/176 (93.2\%), the highest among all four
      evaluated configurations.  14 of 22 tasks received a perfect 8/8.
      No task scored below 5/8, eliminating the outright failure seen
      in the Opus raw run.

\item \textbf{Root cause identification is perfect.}
      The root-cause dimension scored 66/66 (100\%) --- the only run
      across all configurations to achieve this.  The Opus raw run scored
      64/66 (97.0\%), losing 2 points on task~19 where it returned a
      filename with no diagnosis.  Sonnet correctly identifies the root
      cause on every task, including task~19 where it diagnoses the
      \texttt{UnrecognizedUnit.\_\_eq\_\_} issue, even though its patch
      is incomplete.

\item \textbf{File identification is perfect.}
      All 22 tasks scored 44/44.  Direct filesystem access eliminates
      ambiguity about file locations, consistent with all other run
      configurations.

\item \textbf{Patch quality matches Opus raw exactly.}
      The patch dimension scored 54/66 (81.8\%), identical to the Opus
      raw run.  The failure modes differ: Sonnet loses points on
      tasks~8 (\texttt{14096}, over-engineered MRO walk), 13
      (\texttt{14508}, incorrect fallback in \texttt{\_format\_float}),
      and~19 (\texttt{7606}, wrong return value and missing class).
      Opus raw loses points on tasks~7 (\texttt{13977}, partial body
      wrap), 21 (\texttt{8707}, description-only), and~19 (\texttt{7606},
      no patch at all).
      Notably, Sonnet correctly solves task~7 (\texttt{13977}) where
      both Opus raw and all MCP runs scored 6/8.

\item \textbf{Task 7 (\texttt{13977}) is uniquely solved by Sonnet raw.}
      The \texttt{Quantity.\_\_array\_ufunc\_\_} duck-type fix requires
      wrapping the \emph{entire} method body in a \texttt{try/except}
      block and checking \texttt{\_\_array\_ufunc\_\_} on all
      inputs/outputs before returning \texttt{NotImplemented}.  Sonnet
      raw produces this exact structure, scoring 8/8.  Opus raw, Opus
      MCP, and Sonnet MCP all wrap only the input-conversion loop,
      scoring 6/8.  This is the single data point that differentiates
      the overall scores.

\item \textbf{Sonnet is 27\% cheaper than Opus raw.}
      Average cost per task is \$0.54 (Sonnet raw) versus \$0.73
      (Opus raw), a 26\% saving.  Total cost is \$11.83 versus
      \$16.16.  The saving derives from Sonnet's lower per-token rate
      rather than reduced exploration: cache reads are comparable
      (14.96\,M Sonnet vs.\ 14.79\,M Opus), but Sonnet's output is
      56\% larger (432\,K vs.\ 277\,K tokens) because it writes more
      verbose unified diffs and test code.

\item \textbf{Sonnet is 23\% slower than Opus raw.}
      Average wall-clock time is 310\,s (Sonnet) versus 251\,s (Opus).
      Total evaluation time is 6{,}826\,s (\(\approx\)114\,min) versus
      5{,}528\,s (\(\approx\)92\,min).  The primary driver is task~15
      (\texttt{14598}) which ran for 1{,}637\,s and generated
      107{,}959 output tokens --- more than any single task in the Opus
      run.  Excluding this outlier, the per-task time distributions
      are broadly comparable.

\item \textbf{Task 19 (\texttt{7606}) is the only non-trivial failure.}
      Sonnet scores 5/8: correct root cause (3/3), correct file (2/2),
      but incorrect patch (0/3).  The patch fixes
      \texttt{UnrecognizedUnit.\_\_eq\_\_} by returning \texttt{False}
      instead of \texttt{NotImplemented}, and misses the second required
      fix in \texttt{UnitBase.\_\_eq\_\_}.  This is the weakest result
      in the run but still a meaningful improvement over the Opus raw
      fail (3/8, no patch at all).  The issue description focuses on
      \texttt{UnrecognizedUnit}, misdirecting both models.

\item \textbf{Task 13 (\texttt{14508}) is the only case where Sonnet
      introduces a regression.}
      The \texttt{\_format\_float} fix uses \texttt{str(value)} as the
      primary path but falls back to \texttt{f"\{value:.16G\}"} when
      the string exceeds 20 characters.  This fallback reintroduces the
      original precision-expansion bug for large floats.  The ground
      truth uses \texttt{str(value)} unconditionally --- Python's
      \texttt{str(float)} always gives the shortest round-trip
      representation, so no fallback is needed.

\item \textbf{Haiku is absent from 10 of 22 tasks.}
      On tasks 2, 5, 6, 8, 10, 13, 14, 16, 18, 19 Sonnet ran alone with
      no Haiku involvement.  These tend to be smaller, more focused
      fixes (single-file changes, short diffs).  On the 12 tasks where
      Haiku was active, it handled an average of 17.5 requests per task
      for codebase navigation, while Sonnet averaged 14.3 requests for
      reasoning and generation.

\item \textbf{Output tokens are 56\% higher than Opus raw.}
      Total output is 432{,}443 (Sonnet) versus 276{,}939 (Opus).
      Sonnet produces longer patches --- it tends to include parametrized
      tests, regenerated parse tables, and inline commentary.  This
      verbosity contributes to both the higher token cost and the longer
      wall-clock times but does not correlate with higher accuracy: the
      patch dimension score is identical (81.8\%).

\item \textbf{Shared weakness with Opus: task 8 (\texttt{14096}).}
      Both Opus and Sonnet use an over-engineered MRO-walk approach
      for the \texttt{SkyCoord} property error propagation fix,
      scoring 6--7/8 respectively.  The ground truth is a 2-line
      \texttt{\_\_getattribute\_\_} call.  This is a consistent
      failure mode across all configurations, suggesting the issue
      description does not clearly signal the minimal fix.

\end{enumerate}

% ---------------------------------------------------------------------------
\subsection{Cross-Run Summary}\label{sec:raw-sonnet-crossrun}

Table~\ref{tab:crossrun-raw-sonnet} places the Sonnet raw run in context
alongside all other evaluated configurations.

\begin{table}[h]
\centering
\small
\caption{Cross-run comparison: all four configurations.}\label{tab:crossrun-raw-sonnet}
\setlength{\tabcolsep}{5pt}
\begin{tabular}{lcccccrr}
\toprule
\textbf{Configuration} & \textbf{Score} & \textbf{\%} & \textbf{Exact} & \textbf{Near-perf} & \textbf{Fail} & \textbf{Avg Time (s)} & \textbf{Avg Cost (\$)} \\
\midrule
Opus raw      & 162/176 & 92.0\% & 14 &  5 & 1 & 251 & 0.734 \\
Opus MCP      & 163/176 & 92.6\% & 13 &  6 & 0 & 249 & 0.524 \\
Sonnet MCP    & 158/176 & 89.8\% &  7 & 13 & 0 & 149 & 0.300 \\
Sonnet raw    & \textbf{164/176} & \textbf{93.2\%} & \textbf{14} & 5 & \textbf{0} & 310 & \textbf{0.538} \\
\bottomrule
\end{tabular}
\end{table}

\noindent
Sonnet raw achieves the highest accuracy (93.2\%) and the highest
Exact count (14, tied with Opus raw), while eliminating the sole
outright failure and achieving perfect root-cause identification.
The trade-off is a 23\% wall-clock slowdown versus Opus raw, driven
primarily by more verbose output generation rather than additional
exploration depth.  It is the most cost-effective raw-mode
configuration at \$0.54/task.
